\chapter{Defense-Oriented Scientific Positioning}

This appendix summarizes how the project should be defended in front of the jury. It is not a replacement for the main conclusion, but it provides concise academic arguments.

\section{M.1 If the jury asks: Why is this a Digital Twin?}

\begin{itemize}

\item Because the system receives live data from the physical prototype and uses it to update a virtual zone state.

\item Because the virtual state is not only displayed; it is used for prediction, simulation, recommendation, and scheduling.

\item Because the system can influence the physical prototype through validated or manual pump commands.

\item Because it stores history, predictions, simulations, recommendations, commands, and feedback, enabling traceability.

\end{itemize}

\section{M.2 If the jury asks: What is the novelty?}

\begin{itemize}

\item The novelty is not a new sensor or a new ML algorithm alone.

\item The novelty is the integration of Libelium sensing, XBee communication, Raspberry Pi gateway processing, MQTT messaging, FastAPI orchestration, PostgreSQL persistence, machine-learning inference, Digital Twin state, dashboard workflows, WeMos actuation, and pump control into one coherent prototype.

\item The project bridges research state-of-the-art and engineering implementation.

\item It provides a human-supervised Digital Twin workflow with simulation before action.

\end{itemize}

\section{M.3 If the jury asks: What are the weaknesses?}

\begin{itemize}

\item The prototype was tested in a controlled environment, not over a full agricultural season.

\item ML metrics are dataset-level metrics and need real field validation.

\item Security and authentication must be improved for production.

\item Remote sensing and advanced agronomic calibration are future work.

\item The system currently focuses on one experimental zone.

\end{itemize}

\section{M.4 If the jury asks: Why combine Master 2 and Engineering?}

\begin{itemize}

\item The Master 2 part studies the state of the art and scientific positioning of Digital Twins in agriculture.

\item The Engineering part implements and deploys the complete system.

\item The combined thesis is justified because Digital Twin projects require both research understanding and engineering realization.

\end{itemize}
